\section{Introduction}\label{sec:intro}

The Kourovka Notebook~\cite{kourovka} is a collection of unsolved problems in group theory that has been maintained
since 1965; its 21st edition (2026) contains about 2300 problems that were open when the edition went to press,
together with the identity of the proposers and comments on partial progress. It is a natural benchmark for the
question that motivated the experiment reported here: \emph{what can a current general-purpose AI agent, given a
computer algebra system and two days, do with a large collection of genuinely open research problems?}

The experiment was deliberately simple in design. The agent received one instruction (reproduced verbatim in
Section~\ref{sec:prompt}), a machine with GAP~4.16.1 installed, and a time limit of 48 hours; it was asked to
document everything in a local git repository and warned that its solutions would be reviewed independently. The
human experimenter did not intervene except to restart the agent's goal loop after the initial hour and, at the very
end, to ask for this write-up. Everything below---the choice of problems, the computations, the proofs, the
reports, and this paper---was produced by the agent.

\paragraph{Outcome in one paragraph.}
The agent read the whole notebook (partly through delegated sub-agents), selected problems that looked
attackable by finite computation or by a short argument, and worked on about sixty of them. Five problems received
complete answers with self-contained proofs (Section~\ref{sec:solved}, Appendix~\ref{app:proofs}):
\begin{itemize}
\item \textbf{21.32} (P.~J.~Cameron): it is decidable whether a finite group $G$ is the derived subgroup of some
finite group; if it is, a witness of order at most $|\Aut G|\cdot|Z(G)|^2\cdot(2\exp Z(G))^{2\lceil\log_2|G|\rceil}$
exists (a quasi-polynomial bound), so the problem is in nondeterministic quasi-polynomial time.
\item \textbf{16.60} (D.~MacHale): $T(G)=\sum_{\chi\in\Irr(G)}\chi(1)\ge |\{g : \alpha(g)=g^{-1}\}|$ for every
automorphism $\alpha$ of every finite group $G$.
\item \textbf{2.78} (P.~I.~Trofimov; main question): the sets of ``soluble'' and of ``absolutely simple''
group-theoretic numbers are finite; every $k\ge 86$ is the number of $\mathrm{IE}\bar{n}$-systems of a non-soluble
non-simple group. (The second question of 2.78 and the related Problem 3.57, on the distribution of the
non-soluble and simple numbers, remain open.)
\item \textbf{18.46} (Ant.~A.~Klyachko): the bound $|H|\le 2|G|^2$ for a group $H\supseteq G$ in which every
element of $G$ is a square is sharp: $G=D_8$ needs $|H|=128$.
\item \textbf{13.19} (Yu.~M.~Gorchakov): the answer is negative for every prime $p$; in fact every finite
$p$-group occurs as the quotient $Q/H$ in the configuration of the problem.
\end{itemize}
Seven further problems received substantial partial results (Section~\ref{sec:partial}), among them a
structure theorem and a large exclusion range for Berkovich's Problem 16.14, a general theorem and 267 explicitly
realised factor sizes (286 verified factorisations) for Hooshmand's Problem 20.37, and a proof that every minimal simple group violates the hypothesis
of Monakhov's Problem 19.56. About forty-five more problems were
verified in bounded ranges without finding a counterexample (Section~\ref{sec:verifications}).

\paragraph{What this paper is and is not.}
This is a report on an experiment, written by the system that performed it. It records the setting, the process
(including the mistakes), the mathematical content, and an assessment. It is \emph{not} a claim of priority: the
agent had no reliable way to search the literature during the run (web access existed but paywalled sources could
not be read), and it is entirely possible that some of the ``solved'' problems have known answers, that some proofs
reproduce known arguments, or that a problem's status changed after the notebook was compiled. Where the agent
knowingly reused a known theorem (for instance the twisted Frobenius--Schur theorem of Kawanaka and Matsuyama
in Problem 16.60, or Zenkov's theorem on minimal intersections in Problem 21.26) this is stated. We ask readers who
know the literature to treat the results as candidates for verification, and we have tried to make verification as
easy as possible: every proof is self-contained, every computation is a script in the repository, and every
computational claim states exactly what was computed.

\paragraph{Organisation.}
Section~\ref{sec:setting} describes the setting (the task, the machine, the software, the agent).
Section~\ref{sec:process} describes how the run proceeded. Sections~\ref{sec:solved}--\ref{sec:verifications}
present the results; the complete proofs are in the appendices, where they can be checked independently of the
narrative. Section~\ref{sec:discussion} discusses reliability, limitations and lessons. Appendix~\ref{app:timeline}
gives the commit-level timeline and Appendix~\ref{app:repo} the layout of the repository.
